\pagebreak
% Esta no es una solución seria, es algo que se me ocurrió para improvisar
% el texto.
\section{Solución}
Normalmente cuando se piensa en todo es tipo de brechas de seguridad, es
directamente el uso de un antivirus y en el caso de las PyMES se le
incorporaría la acción de concientización a partir de conferencias para hacer
una gestión de todo lo que se está haciendo y como se refiere
\textcite{Bueno2022}, aplicar una décima parte del presupuesto. Por lo tanto,
se pensó en la idea de crear un software inspirado en los programas antivirus
y con Ada Test incorporado para prevenir fallos catastróficos sin el
conocimiento de los usuarios.

A su vez, con  este software se logra sacar provecho al uso del Machine Learning
para identificar elementos que hacen ligeras perturbaciones en los archivos y
que conyevan a que se puedan generar vulnerabilidades que gneren los efectos
explicados por \textcite{Mosquera2019} y \textcite{Rosli2019}. Esta noción
fue inspirada por los elementos presentados por \textcite{Aboaoja2022} al
intentar realizar un análisis con el malware que evoluciona con el paso del
tiempo y que influye en aspectos que indirectamente están relacionados con la
informática, y ni hablar lo que ocurre a nivel nacional según
\textcite{CanoMartinez2022}.
